\chapter{乘积测度}
\label{chap:product-measure}

多重积分的换序问题先要回答两个集合论问题：哪些平面集合可测，以及它们的
测度怎样由矩形的测度决定。若只把“先对一个变量积分”当作计算规则，就会
遗漏截面的可测性，也无法说明何时可以交换次序。本章在
\(\sigma\)-有限测度空间上构造乘积测度，并分别处理非负函数和可积函数。

\section{乘积 \texorpdfstring{\(\sigma\)}{sigma}-代数}

设 \((X,\mathcal A)\)、\((Y,\mathcal B)\) 为可测空间。由可测矩形
\[
 A\times B,\qquad A\in\mathcal A,\ B\in\mathcal B
\]
生成的 \(\sigma\)-代数称为\term{乘积 \(\sigma\)-代数}，记作
\[
 \mathcal A\otimes\mathcal B
 =\sigma\{A\times B:A\in\mathcal A,\ B\in\mathcal B\}.
\]
坐标投影在乘积可测空间上可测。若 \(f:X\to X'\)、\(g:Y\to Y'\) 可测，
则 \((x,y)\mapsto(f(x),g(y))\) 可测。

在拓扑空间上必须区分
\(\mathcal B(X)\otimes\mathcal B(Y)\) 与
\(\mathcal B(X\times Y)\)。对第二可数空间（特别是欧氏空间），二者相等：
开矩形构成可数基，每个开集是可数个开矩形之并。缺少可数性条件时，等式
可能失败。

\section{矩形上的预测度}

设 \((X,\mathcal A,\mu)\)、\((Y,\mathcal B,\nu)\) 为
\(\sigma\)-有限测度空间。希望乘积测度满足
\[
 (\mu\times\nu)(A\times B)=\mu(A)\nu(B).
\]
右端在一个因子为零、另一个为无穷时按 \(0\cdot\infty=0\) 解释。

有限个可测矩形的并可按各坐标集合及其补集作共同细分，写成有限个两两
不交可测矩形的并。这些集合对有限并和差封闭，并包含 \(X\times Y\)，
因而组成集合代数。以下仍称它为矩形环。在其上定义
\[
 \pi\left(\bigsqcup_{k=1}^{m}A_k\times B_k\right)
 =\sum_{k=1}^{m}\mu(A_k)\nu(B_k).
\]

\begin{proposition}
\label{prop:rectangle-ring-premeasure}
上述定义与不交矩形表示无关，并且 \(\pi\) 是该集合环上的预测度。
\end{proposition}

\begin{proof}
记 \(E_x=\{y\in Y:(x,y)\in E\}\) 为集合的竖截面。
设 \(E=\bigsqcup_{k=1}^{m}A_k\times B_k\)。对固定的 \(x\)，所有满足
\(x\in A_k\) 的 \(B_k\) 两两不交，因此
\[
 \nu(E_x)=\sum_{k=1}^{m}\mathbf1_{A_k}(x)\nu(B_k).
\]
右侧是非负扩充实值可测函数；当某个 \(\nu(B_k)=\infty\) 时，
它未必是按本书定义取有限值的简单函数。先把每个系数替换为
\(\min\{\nu(B_k),N\}\)，再让 \(N\uparrow\infty\)，由单调收敛得到
\[
 \int_X\nu(E_x)\,d\mu(x)
 =\sum_{k=1}^{m}\mu(A_k)\nu(B_k).
\]
左侧只由集合 \(E\) 决定，故 \(\pi(E)\) 与所选表示无关；有限可加性也由
截面对不交并保持不交立即得到。

再设 \(E_n\) 是环中两两不交的集合，且
\(E=\bigsqcup_{n=1}^{\infty}E_n\) 仍属于该环。对每个 \(x\)，截面
\((E_n)_x\) 两两不交，所以
\[
 \nu(E_x)=\sum_{n=1}^{\infty}\nu((E_n)_x).
\]
每个截面测度函数都非负可测。对右侧部分和使用单调收敛定理，得到
\[
 \pi(E)=\int_X\nu(E_x)\,d\mu
 =\sum_{n=1}^{\infty}\int_X\nu((E_n)_x)\,d\mu
 =\sum_{n=1}^{\infty}\pi(E_n).
\]
这正是预测度的可数可加性。
\end{proof}

第 \ref{chap:outer-measures} 章的 Carathéodory 延拓定理于是给出
\(\mathcal A\otimes\mathcal B\) 上的测度。

\begin{theorem}[乘积测度]
\label{thm:product-measure-existence}
若 \((X,\mathcal A,\mu)\)、\((Y,\mathcal B,\nu)\) 都是
\(\sigma\)-有限测度空间，则在 \(\mathcal A\otimes\mathcal B\) 上存在
唯一测度 \(\mu\times\nu\) 满足
\[
 (\mu\times\nu)(A\times B)=\mu(A)\nu(B)
\]
对所有可测矩形成立。
\end{theorem}

\begin{proof}
上述预测度经 Carathéodory 定理延拓到矩形生成的
\(\sigma\)-代数。因 \(X=\cup_iX_i\)、\(Y=\cup_jY_j\)，其中各
\(\mu(X_i),\nu(Y_j)<\infty\)，故
\(X\times Y=\cup_{i,j}X_i\times Y_j\)，预测度是
\(\sigma\)-有限的。延拓唯一性定理给出唯一性。
\end{proof}

\section{截面}

对 \(E\subset X\times Y\)，定义
\[
 E_x=\{y\in Y:(x,y)\in E\},\qquad
 E^y=\{x\in X:(x,y)\in E\}.
\]
对函数 \(f:X\times Y\to\overline{\mathbb R}\)，记
\(f_x(y)=f(x,y)\)、\(f^y(x)=f(x,y)\)。

\begin{proposition}[截面的可测性]
\label{prop:measurable-sections}
若 \(E\in\mathcal A\otimes\mathcal B\)，则每个 \(E_x\in\mathcal B\)、
\(E^y\in\mathcal A\)。若 \(f\) 是乘积空间上的扩充实值可测函数，则每个
\(f_x,f^y\) 可测。
\end{proposition}

\begin{proof}
令
\[
 \mathcal C=\{E\subset X\times Y:E_x\in\mathcal B
 \text{ 对所有 }x\}.
\]
\(\mathcal C\) 是 \(\sigma\)-代数，且包含所有可测矩形，故包含
\(\mathcal A\otimes\mathcal B\)。另一方向同理。函数结论由
\(\{y:f_x(y)>a\}=(\{f>a\})_x\) 和半直线判据得到。
\end{proof}

截面可测不意味着任意截面的测度函数自动可测；后者是 Tonelli 定理证明中
需要建立的内容。

\section{集合形式的 Cavalieri 公式}

\begin{theorem}
\label{thm:cavalieri-sets}
设 \((X,\mathcal A,\mu)\)、\((Y,\mathcal B,\nu)\) 为
\(\sigma\)-有限测度空间。若 \(E\in\mathcal A\otimes\mathcal B\)，则
\[
 x\longmapsto\nu(E_x),\qquad y\longmapsto\mu(E^y)
\]
分别可测，并且
\[
 (\mu\times\nu)(E)
 =\int_X\nu(E_x)\,d\mu(x)
 =\int_Y\mu(E^y)\,d\nu(y).
\]
\end{theorem}

\begin{proof}
先设 \(\mu(X),\nu(Y)<\infty\)。令 \(\mathcal D\) 为使第一项可测且
第一等式成立的集合 \(E\)。矩形属于 \(\mathcal D\)，因为
\(\nu((A\times B)_x)=\mathbf1_A(x)\nu(B)\)。若 \(E\in\mathcal D\)，
则
\[
 \nu((E^c)_x)=\nu(Y)-\nu(E_x),
\]
故补集仍在 \(\mathcal D\)。对两两不交的 \(E_n\in\mathcal D\)，截面
也两两不交，测度的可数可加性和单调收敛定理给出
\[
 \nu\bigl((\cup_nE_n)_x\bigr)=\sum_n\nu((E_n)_x)
\]
及积分公式。因此 \(\mathcal D\) 是包含矩形的 Dynkin 系；矩形构成
\(\pi\)-系，由 \(\pi\)-\(\lambda\) 定理，
\(\mathcal D=\mathcal A\otimes\mathcal B\)。

一般 \(\sigma\)-有限情形取递增可测集
\(X_i\uparrow X\)、\(Y_j\uparrow Y\)，其中
\(\mu(X_i),\nu(Y_j)<\infty\)。固定 \(j\)。对每个 \(i\)，把有限情形应用于
\(X_i\times Y_j\) 中的集合 \(E\cap(X_i\times Y_j)\)，得到
\[
 (\mu\times\nu)\bigl(E\cap(X_i\times Y_j)\bigr)
 =\int_{X_i}\nu(E_x\cap Y_j)\,d\mu(x).
\]
这里乘积测度限制在 \(X_i\times Y_j\) 上，与两个限制测度的有限乘积相同；
这可由它们在矩形上一致及延拓唯一性得到。把右端的被积函数在
\(X\setminus X_i\) 上补为零。随 \(i\) 增大，这些函数逐点递增到
\(x\mapsto\nu(E_x\cap Y_j)\)。因此测度的上升连续性和单调收敛定理给出
\[
 (\mu\times\nu)\bigl(E\cap(X\times Y_j)\bigr)
 =\int_X\nu(E_x\cap Y_j)\,d\mu(x),
\]
并同时说明右端的被积函数可测。最后令 \(j\to\infty\)。左侧递增到
\((\mu\times\nu)(E)\)，而
\(\nu(E_x\cap Y_j)\uparrow\nu(E_x)\)；再用一次单调收敛定理，得到第一个
等式及 \(x\mapsto\nu(E_x)\) 的可测性。交换 \(X,Y\) 即得第二个等式。
\end{proof}

\section{Tonelli 定理}

\begin{theorem}[Tonelli]
\label{thm:tonelli}
设 \((X,\mathcal A,\mu)\)、\((Y,\mathcal B,\nu)\) 为
\(\sigma\)-有限测度空间。若 \(f:X\times Y\to[0,+\infty]\)
对乘积 \(\sigma\)-代数可测，则
\[
 x\longmapsto\int_Yf(x,y)\,d\nu(y)
\]
可测，并且
\[
 \int_{X\times Y}f\,d(\mu\times\nu)
 =\int_X\left(\int_Yf(x,y)\,d\nu(y)\right)d\mu(x)
 =\int_Y\left(\int_Xf(x,y)\,d\mu(x)\right)d\nu(y).
\]
三者可以同时为 \(+\infty\)。
\end{theorem}

\begin{proof}
对指示函数，这是定理 \ref{thm:cavalieri-sets}。由非负简单函数的线性，
结论对非负简单函数成立。对一般 \(f\)，取
\(\varphi_n\uparrow f\)。对每个 \(x\)，单调收敛定理给出
\[
 \int_Y\varphi_n(x,y)\,d\nu\uparrow\int_Yf(x,y)\,d\nu.
\]
左侧是可测函数，所以极限可测。再次分别在 \(X\) 和 \(X\times Y\) 上
使用单调收敛定理，得到第一种迭代积分公式；另一种次序同理。
\end{proof}

\section{Fubini 定理}

\begin{theorem}[Fubini]
\label{thm:fubini}
设 \((X,\mathcal A,\mu)\)、\((Y,\mathcal B,\nu)\) 为
\(\sigma\)-有限测度空间，\(f\) 是乘积空间上的可测实函数，且
\(\int|f|\,d(\mu\times\nu)<\infty\)。则
\[
 f_x\in L^1(\nu)\quad\text{对几乎所有 }x,
\qquad
 f^y\in L^1(\mu)\quad\text{对几乎所有 }y.
\]
在内层绝对值积分为无穷的可测零集上，把相应的内层积分定义为零，
所得函数可积，且
\[
 \int_{X\times Y}f\,d(\mu\times\nu)
 =\int_X\int_Yf(x,y)\,d\nu(y)\,d\mu(x)
 =\int_Y\int_Xf(x,y)\,d\mu(x)\,d\nu(y).
\]
\end{theorem}

\begin{proof}
对 \(|f|\) 使用 Tonelli 定理：
\[
 \int_X\left(\int_Y|f_x|\,d\nu\right)d\mu
 =\int_{X\times Y}|f|\,d(\mu\times\nu)<\infty.
\]
令 \(N_X=\{x:\int_Y|f_x|\,d\nu=\infty\}\)。由 Tonelli 定理，
它可测且 \(\mu(N_X)=0\)。在补集上，两个可测函数
\(\int_Y f_x^+\,d\nu\)、\(\int_Y f_x^-\,d\nu\) 都有限，故其差可测；
在 \(N_X\) 上置零仍可测。其绝对值不超过内层绝对值积分，因此可积。
分别对 \(f^+,f^-\) 使用 Tonelli 再相减，得到第一个迭代公式；
交换变量得到第二个。这一处理不要求因子测度完备。
\end{proof}

复函数的积分按实部、虚部分别定义。若 \(\int|f|<\infty\)，两个分量
均可积，分别应用上述实函数版本即得相同的迭代积分公式。

实用上常先对 \(|f|\) 作一次迭代积分。若结果有限，Fubini 的全部结论随即
成立。若只算带符号的内层积分并得到有限值，仍不足以交换次序。

\section{完备化与例外截面}

以上定理首先建立在 \(\mathcal A\otimes\mathcal B\) 上。即使两个因子测度
都是完备的，乘积测度也不一定完备。对乘积测度作完备化以后，
“每个截面可测”应相应改为“几乎每个截面可测”。

设两个因子测度完备，\(E\) 属于乘积测度的完备化。按完备化的定义，可取
\(B,Z\in\mathcal A\otimes\mathcal B\)，使
\[
 E\mathbin{\triangle}B\subset Z,\qquad(\mu\times\nu)(Z)=0.
\]
由 Cavalieri 公式，\(\nu(Z_x)=0\) 对几乎所有 \(x\) 成立。在这些点，
\(E_x\mathbin{\triangle}B_x\subset Z_x\)，因子测度的完备性保证
\(E_x\) 可测，且 \(\nu(E_x)=\nu(B_x)\)。故截面积分公式仍成立，但内层
积分只需在几乎所有外层参数上定义。

函数形式也如此。任一非负完备化可测函数 \(f\) 都有一个
\(\mathcal A\otimes\mathcal B\)-可测函数 \(g\)，在某个乘积零测集以外与它相等。
为证明这一点，取非负简单函数 \(s_n\uparrow f\)，把各 \(s_n\) 中有限个
完备化可测集换成与其仅相差零测集的乘积可测集。所有差异包含在可数个
乘积零测集的并 \(Z\) 中。所得函数序列在 \(Z\) 外仍趋于 \(f\)，
令 \(g\) 在 \(Z\) 外为其上极限，在 \(Z\) 上为零，即得所求。
对 \(g\) 应用 Tonelli 定理，再对 \(Z\) 应用截面零测性，便得到
完备化版本。可积实函数拆成正负部，复函数再拆成实部、虚部，Fubini 公式
因而也保留。

欧氏空间中的 Lebesgue 测度正是 Borel 测度的完备化；Borel 乘积测度在
矩形上给通常的体积，唯一性保证二者相符。因此上述版本适用于一般
Lebesgue 可测函数，而不只适用于 Borel 可测函数。
例外截面不能删去：若 \(V\subset[0,1]\) 不 Lebesgue 可测，则
\(\{0\}\times V\) 包含在平面 Borel 零测集 \(\{0\}\times[0,1]\) 中，
所以平面 Lebesgue 可测；但在 \(x=0\) 处的截面恰为 \(V\)。

\section{二重级数与反例}

在 \(X=Y=\mathbb N\) 上取计数测度，Tonelli 与 Fubini 分别成为：
\[
 a_{mn}\ge0\quad\Longrightarrow\quad
 \sum_m\sum_na_{mn}=\sum_n\sum_ma_{mn},
\]
以及
\[
 \sum_{m,n}|a_{mn}|<\infty
 \quad\Longrightarrow\quad
 \sum_m\sum_na_{mn}=\sum_n\sum_ma_{mn}.
\]

\begin{counterexample}
令
\[
 a_{mn}=
 \begin{cases}
 1,&n=m,\\
 -1,&n=m+1,\\
 0,&\text{其他}.
 \end{cases}
\]
每一行的和为零，所以 \(\sum_m(\sum_na_{mn})=0\)。第 \(0\) 列的和为
\(1\)，其余各列的和为零，所以
\(\sum_n(\sum_ma_{mn})=1\)。绝对值二重和发散，Fubini 条件不成立。
\end{counterexample}

\begin{counterexample}[缺少 \(\sigma\)-有限性]
在 \([0,1]\) 的 Borel 集上，第一坐标取 Lebesgue 测度 \(m\)，第二坐标
取计数测度 \(c\)。后者不是 \(\sigma\)-有限的，因为每个有限测度集只有
有限多个点。对角线 \(D=\{(x,y):x=y\}\) 是 Borel 集，但
\[
 \int_{[0,1]}c(D_x)\,dm(x)=\int_{[0,1]}1\,dm=1,
 \qquad
 \int_{[0,1]}m(D^y)\,dc(y)=\int_{[0,1]}0\,dc=0.
\]
所以即使函数非负且乘积可测，两种迭代积分也未必相等。这个例子否定的是
无条件换序，并不声称无法构造任何满足矩形公式的乘积测度。
\end{counterexample}

\section*{习题}
\addcontentsline{toc}{section}{习题}
\subsection*{A组　乘积可测结构}
\begin{enumerate}[label=\textbf{\thechapter.\arabic*.}]
\exerciseitem{product-sigma-coordinate}
设 \(f:X\to X'\)、\(g:Y\to Y'\) 为可测空间之间的可测映射。
证明两个坐标投影可测，并证明它们组成的
\((x,y)\mapsto(f(x),g(y))\) 可测。

\exerciseitem{borel-product-euclidean}
证明
\(\mathcal B(\mathbb R^m)\otimes\mathcal B(\mathbb R^n)
=\mathcal B(\mathbb R^{m+n})\)。

\exerciseitem{product-rectangle-refinement}
证明有限个可测矩形的并可以写成有限个两两不交可测矩形的并，并据此验证
矩形环上预测度定义的良定义性。

\exerciseitem{section-algebra}
验证 \((E^c)_x=(E_x)^c\) 和
\((\cup_nE_n)_x=\cup_n(E_n)_x\)，补全截面可测性的证明。

\end{enumerate}
\subsection*{B组　Tonelli 与 Fubini}
\begin{enumerate}[label=\textbf{\thechapter.\arabic*.},resume]
\exerciseitem{cavalieri-graph}
设 \(f:\mathbb R\to\mathbb R\) Borel 可测。证明其图像
\(\Gamma_f=\{(x,f(x)):x\in\mathbb R\}\) 是 Borel 集，并用截面公式证明
\(m_2(\Gamma_f)=0\)。

\exerciseitem{tonelli-layer-cake}
用 Tonelli 定理证明对 \(f\ge0\)，
\[
 \int_Xf\,d\mu=\int_0^\infty\mu(\{f>t\})\,dt.
\]

\exerciseitem{tonelli-region-integral}
利用截面计算集合
\[
 E=\{(x,y)\in(0,\infty)^2:x+y<1\}
\]
的二维 Lebesgue 测度，并交换积分次序复核。

\exerciseitem{fubini-section-null}
若 \(E\subset X\times Y\) 可测且
\((\mu\times\nu)(E)=0\)，证明 \(E_x\) 对几乎所有 \(x\) 为
\(\nu\)-零测集。说明“每个截面都零测”是否必然推出 \(E\) 零测。

\end{enumerate}
\subsection*{C组　换序、卷积与推广}
\begin{enumerate}[label=\textbf{\thechapter.\arabic*.},resume]
\exerciseitem{fubini-integrability-test}
讨论
\[
 f(x,y)=\frac{x^2-y^2}{(x^2+y^2)^2}
\]
在 \((-1,1)^2\) 上的绝对可积性（原点处补为零），并比较两种先后次序
的反常迭代积分；外层参数为零时的发散截面应单独说明。

\exerciseitem{double-series-counterexample}
逐项验证正文二重级数反例的两种迭代和，并证明其绝对值二重和发散。

\exerciseitem{completed-product-exceptional-section}
取一个不 Lebesgue 可测的 \(V\subset[0,1]\)。证明
\(\{0\}\times V\) 属于二维 Lebesgue \(\sigma\)-代数，却不属于
\(\mathcal L(\R)\otimes\mathcal L(\R)\)。据此说明两个完备测度的乘积
不一定完备，并指出 Fubini 定理中为什么需要允许例外截面。

\optionalexerciseitem{convolution-measurability}
若 \(f,g:\mathbb R^n\to[0,\infty]\) 可测，证明
\[
 (f*g)(x)=\int_{\mathbb R^n}f(x-y)g(y)\,dy
\]
是可测函数，并证明
\(\int(f*g)=(\int f)(\int g)\)，允许无穷值。

\projectexerciseitem{associativity-product-measure}
对三个 \(\sigma\)-有限测度空间，证明
\((\mu\times\nu)\times\lambda\) 与
\(\mu\times(\nu\times\lambda)\) 在自然识别下相同，并给出三重 Tonelli
和 Fubini 公式。
\end{enumerate}
